L'experiment d'Oersted fet el 1820 (figura 1) consisteix en un fil conductor paral·lel a la component horitzontal del camp magnètic terrestre per on passa un corrent elèctric i una agulla magnètica just a sobre o a sota del fil. Oersted observà que l'orientació de l'agulla canviava si el fil connectat estava situat a sobre o a sota de l'agulla.

\figura[0.35]{fig1.png}
{\centering\small \textsc{Figura 1.} Experiment d'Oersted.\par}\medskip

En la figura 2 no sabem si el fil conductor està situat per sobre o per sota de la brúixola.

\figura[0.45]{fig2.png}
{\centering\small \textsc{Figura 2.} \textit{a}) Circuit obert. \textit{b}) Circuit tancat.\par}\medskip

\begin{apartat}{1,25}
Representeu sobre l'agulla de la brúixola del dibuix \textit{b} de la figura 2 els vectors del camp magnètic terrestre ($B_\mathrm{Terra}$) i del camp magnètic generat pel fil ($B_\mathrm{fil}$). Argumenteu si la brúixola del dibuix \textit{b} de la figura 2 està situada a sobre o a sota del fil conductor en aquest cas.
\end{apartat}

\begin{apartat}{1,25}
En aquest muntatge experimental, la separació entre la brúixola i el fil és de 10~cm. Observem que l'agulla de la brúixola forma un angle de $45^\circ$ amb la direcció del corrent elèctric quan circulen 20~A pel fil conductor. Calculeu la component horitzontal del camp magnètic terrestre ($B_\mathrm{Terra}$).
\end{apartat}

\dades{%
  El mòdul del camp magnètic creat per un fil infinit per on circula un corrent $I$ a una distància $r$ del fil és $B = \dfrac{\mu_0\, I}{2\pi\, r}$.\\
  $\mu_0 = 4\pi\times 10^{-7}\ \mathrm{T\,m\,A^{-1}}$.}
